\chapter{The Pump Solve: Harmonic Balance}
\label{chap:pumphb}

This chapter specifies the nonlinear pump solve completely. Every nested loop
is developed from the governing equation, and every formula is given in the
form the implementation actually uses. Where the implementation departs from the
idealised statement --- and there are three such places --- the departure is
recorded in an \emph{Implementation} box with the coded form, its consequence,
and the numerical evidence.

\section{The harmonic-balance formulation}

\subsection{Ansatz}

The pump drives the circuit at a single frequency $\wp$, and we seek the
\emph{periodic steady state} of period $T = 2\pi/\wp$. Expand the node fluxes in
a finite set of positive harmonics $\modeset \subset \mathbb{Z}_{>0}$:

\begin{convention}[positive-phasor reconstruction]
\label{conv:phasor}
\begin{equation}
  \boxed{\;
  \xvec(t) = 2\,\Real\!\left[\sum_{k\in\modeset} X_k\,e^{+ik\wp t}\right],
  \qquad X_k\in\mathbb{C}^{n}.
  \;}
  \label{eq:synthesis}
\end{equation}
The unknown is the stacked array $\Xvec = \{X_k\}_{k\in\modeset}\in
\mathbb{C}^{H\times n}$ with $H = |\modeset|$.
\end{convention}

The factor $2$ and the sign $+i$ are the entire content of the convention. Under
it, a single mode $X_1 = \tfrac12 A$ reconstructs to $A\cos(\wp t)$, so the
phasor is \emph{half} the physical amplitude. Every source coefficient, every
power conversion, and every imported pump solution must respect this or the
result is wrong by factors of two.

\subsection{The residual}

Substituting \eqref{eq:synthesis} into the governing
equation~\eqref{eq:governing} and projecting onto $e^{-ik\wp t}$ yields, for
each retained mode $k$,
\begin{equation}
  \boxed{\;
  R_k(\Xvec;\lambda)
  \;=\;
  \Dmat_k X_k + N_k(\Xvec) - \lambda \Svec_k
  \;=\; 0,
  \;}
  \label{eq:residual-k}
\end{equation}
where $\Dmat_k$ is the linear dynamic block (\cref{sec:freq-loop}), $N_k$ the
projected nonlinear branch current (\cref{sec:residual-eval}), $\Svec_k$ the
source coefficient (\cref{sec:power-loop}), and $\lambda\in[0,1]$ the
continuation parameter (\cref{sec:continuation}).

Stacking over modes,
\begin{equation}
  \Rvec(\Xvec;\lambda)
  = \mathcal{D}\Xvec + \Nvec(\Xvec) - \lambda\Svec
  = \vect{0},
  \qquad
  \mathcal{D} = \blockdiag(\Dmat_{k_1},\dots,\Dmat_{k_H}).
  \label{eq:residual-stacked}
\end{equation}

This is the equation every layer below works toward solving. It is a nonlinear
algebraic system in $2Hn$ real unknowns (real and imaginary parts of every
phasor).

\subsection{Loop hierarchy}

The solve is organised as eight nested layers. Each answers one question:

\begin{center}
\begin{tabular}{@{}r l l@{}}
\toprule
& \textbf{Layer} & \textbf{Question} \\
\midrule
1 & Frequency      & Which pump frequency? \\
2 & Power          & Which pump amplitude? \\
3 & Continuation   & Can we reach the full drive gradually? \\
4 & Newton         & Where is the nonlinear root? \\
5 & \quad Residual & How wrong is the current guess? \\
6 & \quad Preconditioner & What cheap $\Mmat\approx\Jmat$ can we build? \\
7 & \quad Krylov (GMRES + JVP) & What correction solves the linearisation? \\
8 & \quad Line search & How much of that correction is safe? \\
\bottomrule
\end{tabular}
\end{center}

\Cref{sec:hierarchy} collects the whole thing in one place.

\section{Layer 1: the frequency loop}
\label{sec:freq-loop}

\subsection{The linear dynamic block}

For a harmonic $k$, the ansatz \eqref{eq:synthesis} gives
$\dot{\xvec}_k = ik\wp X_k$ and $\ddot{\xvec}_k = -(k\wp)^{2}X_k$, so the linear
part of \eqref{eq:governing} contributes
\begin{equation}
  \Cmat\ddot{\xvec}_k + \Gmat\dot{\xvec}_k + \Kmat\xvec_k
  = \Dmat_k X_k
\end{equation}
with
\begin{equation}
  \boxed{\;
  \Dmat_k(\wp) \;=\; \Kmat \;-\; (k\wp)^{2}\,\Cmat \;+\; i\,k\wp\,\Gmat .
  \;}
  \label{eq:dynamic-block}
\end{equation}

Selecting a pump frequency therefore changes the entire family
$\{\Dmat_k\}_{k\in\modeset}$ and hence produces a genuinely different
harmonic-balance problem. The blocks are constant in $\Xvec$, so they are built
once per frequency and reused beneath every inner loop.

\begin{implnote}[\code{\_build\_linear\_blocks}]
\begin{lstlisting}
for k in self.grid.k:
    wk = float(k) * self.grid.omega
    Dk = Kc + (-wk * wk) * Cc + (1j * wk) * Gc
    blocks.append(Dk.tocsc())
\end{lstlisting}
Exactly \eqref{eq:dynamic-block}. Verified against a directly assembled
reference at $k\in\{1,3,5\}$: maximum relative deviation
$\num{0.000e+00}$.
\end{implnote}

\subsection{The pump-mode basis}
\label{sec:pump-basis}

Which modes to retain is a physics decision, not a numerical one.

For an \emph{unbiased} four-wave-mixing device the branch law is odd, so the
pump waveform contains only odd harmonics: $X_k \equiv 0$ for even $k$. Using a
dense basis $\{1,2,\dots,H\}$ therefore spends half the unknowns on modes that
are identically zero while truncating the high odd harmonics that carry real
amplitude. The correct basis is
\begin{equation}
  \modeset_{\mathrm{odd}} = \{1,3,5,\dots,2K-1\}.
\end{equation}

\begin{measured}[basis choice matters]
Replacing the legacy dense basis $\{1,\dots,H\}$ with the odd basis reduced the
gain discrepancy on a reference JTWPA from $\approx\SI{0.89}{\dB}$ RMS to
$\approx\SI{0.0006}{\dB}$ RMS. The dense basis was truncating the high odd pump
content.
\end{measured}

\begin{implnote}[\code{pump/basis.py} as single source of truth]
\code{resolve\_pump\_basis} returns a \code{PumpBasis} under one of four
policies:
\begin{center}
\begin{tabular}{@{}l l l@{}}
\toprule
Policy & Modes & Use for \\
\midrule
\code{positive\_odd\_jc} & $\{1,3,\dots,2K-1\}$ & unbiased 4WM \\
\code{dense\_real} & $\{1,2,\dots,H\}$ & biased / DC / 3WM \\
\code{positive\_phasor\_explicit} & user list & anything else \\
\code{auto\_jc} & inferred & single-pump cases \\
\bottomrule
\end{tabular}
\end{center}
\code{PumpBasis.to\_metadata} writes \code{pump\_modes},
\code{pump\_basis="positive\_phasor"}, \code{real\_reconstruction\_factor=2},
\code{omega\_p} and \code{phase\_convention="exp\_plus\_i\_k\_omega\_t"} into the
saved solution; \code{load\_pump\_basis\_from\_solution} reads them back. A gain
solve therefore cannot silently disagree with the pump solve about what the
stored phasors mean.
\end{implnote}

\subsection{The time grid}

The nonlinearity is evaluated by alternating frequency--time (AFT): synthesise
on a time grid, apply the branch law pointwise, transform back. The grid is
\begin{equation}
  \theta_r = \frac{2\pi r}{N_t},\qquad r = 0,\dots,N_t-1,
  \qquad t_r = \frac{\theta_r}{\wp}.
\end{equation}

\begin{proposition}[alias-free sampling]
To represent products up to $2\max(\modeset)$ without aliasing onto a retained
mode, it is necessary that
\begin{equation}
  N_t \;\ge\; 2\max(\modeset) + 1 .
  \label{eq:nt-bound}
\end{equation}
\end{proposition}

\begin{implnote}[\code{HarmonicGrid}]
\code{\_\_post\_init\_\_} enforces \eqref{eq:nt-bound} and additionally requires
$N_t$ even. It precomputes the synthesis and analysis matrices
\begin{lstlisting}
self.E = np.exp(1j * self.omega * self.t[:, None] * self.k[None, :])
self.E_conj_T_over_nt = self.E.conj().T / self.nt
\end{lstlisting}
so that synthesis is \code{2.0 * np.real(self.E @ X)} and projection is
\code{self.E\_conj\_T\_over\_nt @ y\_t}. Note this is a dense
$(N_t\times H)$ matrix product, not an FFT: $H$ is small (typically 10--20) and
the mode list is not contiguous, so the general form is both simpler and
competitive.
\end{implnote}

\section{Layer 2: the power loop}
\label{sec:power-loop}

\subsection{Power to current}

The pump generator power is converted to an on-chip peak current by
\eqref{eq:dbm-to-w}--\eqref{eq:w-to-i}, then scaled by the convention factor
$s_{\mathrm{JC}}$ of \cref{sec:power-conversion}:
\begin{equation}
  I_{\mathrm{solver}} = s_{\mathrm{JC}}\,I_{\mathrm{peak}},
  \qquad
  I_{\mathrm{peak}} = \sqrt{\frac{2P_{\mathrm{device}}}{Z_0}} .
\end{equation}

\subsection{The source coefficient}

Under \cref{conv:phasor}, a physical drive $I\cos(k_{\mathrm{src}}\wp t)$ at the
pump node corresponds to a phasor of $\tfrac12 I$. Including the continuation
scale,

\begin{equation}
  \boxed{\;
  \Svec_k(\lambda)
  = \tfrac{1}{2}\,\lambda\,I_{\mathrm{solver}}\;
    \vect{e}_{p}\;\delta_{k,k_{\mathrm{src}}},
  \;}
  \label{eq:source-coeff}
\end{equation}
where $\vect{e}_p$ selects the pump node and $k_{\mathrm{src}}$ is the driven
mode (normally the fundamental). The physical source waveform is then
\begin{equation}
  \isrc(t) = 2\,\Real\!\left[\Svec_{k_{\mathrm{src}}}
             e^{ik_{\mathrm{src}}\wp t}\right].
  \label{eq:source-time}
\end{equation}

\begin{implnote}[\code{source\_coeffs}]
\begin{lstlisting}
def source_coeffs(self, source_scale):
    S = np.zeros((self.H, self.n), dtype=np.complex128)
    source_value = 0.5 * source_scale * self.pump_current_a
    S[self.source_row, self.pump_node_index] = source_value
    return S
\end{lstlisting}
Exactly \eqref{eq:source-coeff}, with \code{source\_row} the position of
$k_{\mathrm{src}}$ in the mode list.
\end{implnote}

\begin{implnote}[\code{source\_time} ignores $k_{\mathrm{src}}$ --- a latent defect]
\label{impl:source-time}
The companion time-domain routine does \emph{not} implement
\eqref{eq:source-time}. It hardcodes the fundamental:
\begin{lstlisting}
def source_time(self, source_scale):
    src = np.zeros((self.grid.nt, self.n), dtype=float)
    src[:, self.pump_node_index] = (
        source_scale * self.pump_current_a
        * np.cos(self.grid.omega * self.grid.t))   # <-- no k_src factor
    return src
\end{lstlisting}
Measured relative deviation between $2\Real[\Svec_k e^{ik\wp t}]$ and
\code{source\_time()}:
\begin{center}
\begin{tabular}{@{}r r@{}}
\toprule
\code{source\_mode} & deviation \\
\midrule
1 & $\num{0.000e+00}$ \\
3 & $\num{1.536}$ \\
\bottomrule
\end{tabular}
\end{center}
\code{source\_time} is consumed only by \code{time\_residual}, an optional
diagnostic gated on \code{compute\_time\_residual}. It never enters the Newton
path, so no published result is affected. It is nonetheless wrong for any
$k_{\mathrm{src}}\neq 1$ and is recorded here as a known defect.
\end{implnote}

\subsection{Warm-start prediction across the power axis}

Consecutive power points are close, so the previous solution is an excellent
initial guess and the previous \emph{two} are better still. Given converged
states at injected currents $I_{j-1}, I_j$, predict at $I_{j+1}$ by linear
extrapolation in the natural coordinate:
\begin{equation}
  \boxed{\;
  \Xvec_{\mathrm{pred}}
  = \Xvec_j
  + \frac{I_{j+1}-I_j}{I_j - I_{j-1}}\left(\Xvec_j - \Xvec_{j-1}\right).
  \;}
  \label{eq:power-secant}
\end{equation}

The abscissa is the injected \emph{current}, not the power in \si{\dBm},
because the source enters \eqref{eq:residual-k} linearly in current.

\begin{implnote}[\code{axis\_secant}]
\file{pump/predictors.py} implements \eqref{eq:power-secant} generically over an
abscissa $t$:
\begin{lstlisting}
beta  = (t_target - t_prev1) / (t_prev1 - t_prev2)
guess = x_prev1 + beta * (x_prev1 - x_prev2)
\end{lstlisting}
returning \code{None} on shape mismatch, missing abscissa, or
$|t_{j}-t_{j-1}|<10^{-30}$. The same routine serves the frequency axis with
$t=f_p$. Additional predictors --- \code{copy}, \code{corner}, \code{plane},
and a residual-ranked \code{portfolio} --- are selectable per campaign.
\end{implnote}

\section{Layer 3: the continuation loop}
\label{sec:continuation}

\subsection{Natural-parameter continuation}

A cold Newton solve at full pump power generally diverges: the initial guess is
far outside the basin of attraction. The remedy is to embed the target problem
in a homotopy,
\begin{equation}
  \boxed{\;
  \Rvec(\Xvec;\lambda) = \mathcal{D}\Xvec + \Nvec(\Xvec) - \lambda\Svec = 0,
  \qquad 0\le\lambda\le 1,
  \;}
  \label{eq:homotopy}
\end{equation}
and track the solution branch from $\lambda=0$ (no drive; $\Xvec=\vect{0}$ is an
exact solution) to $\lambda=1$ (the physical problem). Since the source enters
linearly, this is natural-parameter continuation in the source amplitude.

\subsection{Fixed ladder}

The simplest scheme takes $N_\lambda$ equal steps,
\begin{equation}
  \lambda_j = \frac{j}{N_\lambda},\qquad j=1,\dots,N_\lambda,
\end{equation}
solving $\Rvec(\Xvec_j;\lambda_j)=0$ with initial guess
$\Xvec_j^{(0)} = \Xvec_{j-1}$.

\begin{implnote}[resumable fixed ladder]
\code{solve\_continuation} accepts \code{lambda\_start} (default $0.0$). This
exists so the adaptive fallback below can \emph{resume} rather than restart; all
other callers are unaffected.
\end{implnote}

\subsection{Adaptive stepping}

A fixed ladder wastes effort in the easy low-$\lambda$ region and is too coarse
near a fold. The adaptive scheme sets a target
\begin{equation}
  \lambda_{\mathrm{target}} = \min(1,\;\lambda + \Delta\lambda),
\end{equation}
attempts a Newton solve there, and updates the step by
\begin{equation}
  \Delta\lambda \leftarrow
  \begin{cases}
    \min(1,\;g\,\Delta\lambda), & \text{on success},\quad g \approx 1.5,\\[2pt]
    s\,\Delta\lambda,           & \text{on failure},\quad s \approx 0.5 .
  \end{cases}
  \label{eq:adaptive-step}
\end{equation}

\begin{implnote}[step-control constants]
\code{solve\_adaptive\_continuation} takes \code{growth} and \code{shrink} as
arguments; the wrapper that campaigns use fixes them at
\code{growth=1.5}, \code{shrink=0.5}, matching \eqref{eq:adaptive-step}. The
step is additionally clamped by \code{step = min(step, 1.0 - lambda\_current)}
so a target never overshoots $\lambda=1$, and floored at \code{min\_step}.
\end{implnote}

\subsection{Predictors along the branch}

\paragraph{Secant.} With two converged points on the branch,
\begin{equation}
  \Xvec_{\mathrm{pred}}
  = \Xvec_j
  + \frac{\lambda_{\mathrm{target}}-\lambda_j}{\lambda_j-\lambda_{j-1}}
    \left(\Xvec_j - \Xvec_{j-1}\right).
  \label{eq:lambda-secant}
\end{equation}
Structurally identical to \eqref{eq:power-secant} with $t=\lambda$.

\paragraph{Tangent.} Differentiating \eqref{eq:homotopy} along the branch,
\begin{equation}
  \Jmat(\Xvec)\frac{\mathrm{d}\Xvec}{\mathrm{d}\lambda}
  + \frac{\partial\Rvec}{\partial\lambda} = 0,
  \qquad
  \frac{\partial\Rvec}{\partial\lambda} = -\Svec,
\end{equation}
so the branch tangent solves
\begin{equation}
  \boxed{\;\Jmat(\Xvec)\,\dot{\Xvec} = \Svec,\;}
  \qquad
  \Xvec_{\mathrm{pred}} = \Xvec + \Delta\lambda\,\dot{\Xvec}.
  \label{eq:tangent-predictor}
\end{equation}
This is an Euler predictor along the solution curve and costs one extra linear
solve per step.

\begin{implnote}[\code{tangent\_predictor}]
\begin{lstlisting}
S = problem.source_coeffs(1.0) - problem.source_coeffs(0.0)
\end{lstlisting}
The difference form is used rather than \code{source\_coeffs(1.0)} alone, so
that a problem with an affine (rather than purely linear) source path --- as in
the multitone signal substep of \cref{sec:signal-substep} --- gives the correct
$\partial\Rvec/\partial\lambda$ without modification.
\end{implnote}

\subsection{Fallback, and why it resumes}

When adaptive bisection drives $\Delta\lambda$ below \code{min\_step} --- the
signature of a genuine fold in source amplitude, where the reduction ratio
stalls near $1.0$ at every attempt --- the solver falls back to a fixed ladder.

\begin{implnote}[resume, do not restart]
An earlier version passed the \emph{original} seed and $\lambda=0$ to the
fallback, discarding every state the adaptive phase had already converged. On a
real map column this burned the entire per-point deadline getting back only to
$\lambda=0.75$ after adaptive stepping had already reached $\lambda=0.9375$
converged.

The fallback now resumes from $(\Xvec_{\mathrm{current}},
\lambda_{\mathrm{current}})$ and sizes itself to the \emph{remaining} span at
the original granularity:
\begin{lstlisting}
step_size = 1.0 / fallback_fixed_steps
fallback_steps = max(1, math.ceil(remaining / step_size - 1e-9))
\end{lstlisting}
\end{implnote}

\subsection{Advanced continuation}

Two further schemes are available for stiff cases:

\paragraph{Pseudo-arclength.} Near a fold, $\lambda$ ceases to parameterise the
branch --- the curve turns back. Pseudo-arclength continuation augments the
system with a normalisation constraint and solves for $(\Xvec,\lambda)$
jointly via a bordered linear system, allowing the branch to be followed
around the turning point. \code{solve\_arclength} implements this with a
modified-Newton inner solve.

\paragraph{Pseudo-transient.} \code{solve\_pseudo\_transient} solves
\begin{equation}
  \frac{\Xvec^{(n+1)} - \Xvec^{(n)}}{\delta} + \Rvec(\Xvec^{(n+1)}) = 0
\end{equation}
with a pseudo-timestep $\delta$ grown inversely with the residual drop (switched
evolution relaxation). This damps very stiff starts at the cost of linear rather
than quadratic asymptotic convergence.

\begin{caveat}[experimental status]
Arclength and fold-following are functional but experimental on the stiffest
device studied: fold-following may report no fold in range, and the arclength
target endpoint is linearly interpolated, so it is used as a warm guess rather
than a polished root.
\end{caveat}

\subsection{When continuation is skipped}

Warm-started cells --- the common case in a map, where a neighbouring point has
already converged --- skip continuation entirely and solve directly at
$\lambda=1$. Continuation is for cold, seeded, and recovery solves.

\section{Layer 4: the Newton loop}

\subsection{The correction equation}

At Newton iteration $n$, expand the residual about the current iterate:
\begin{equation}
  \Rvec(\Xvec^{(n)} + \Delta\Xvec)
  \approx
  \Rvec(\Xvec^{(n)}) + \Jmat(\Xvec^{(n)})\,\Delta\Xvec .
\end{equation}
Setting the linearisation to zero gives
\begin{equation}
  \boxed{\;
  \Jmat(\Xvec^{(n)})\,\Delta\Xvec^{(n)} = -\Rvec(\Xvec^{(n)}),
  \;}
  \label{eq:newton}
\end{equation}
and the update, with an accepted damping factor $\alpha_n$
(\cref{sec:line-search}),
\begin{equation}
  \boxed{\;
  \Xvec^{(n+1)} = \Xvec^{(n)} + \alpha_n\,\Delta\Xvec^{(n)} .
  \;}
  \label{eq:newton-update}
\end{equation}

\subsection{Real packing}

$\Rvec$ is \emph{not} a complex-analytic function of $\Xvec$. The nonlinearity
acts on the real waveform $\xvec(t)$, which by \eqref{eq:synthesis} depends on
both $X_k$ and $\overline{X_k}$. There is therefore no complex Jacobian; the
derivative exists only in real coordinates.

\begin{convention}[real packing]
\begin{equation}
  \mathrm{pack}(\Xvec) =
  \begin{bmatrix} \Real\Xvec \\ \Imag\Xvec \end{bmatrix}
  \in\mathbb{R}^{2Hn},
\end{equation}
with the real part of the entire $(H,n)$ array first, then the imaginary part.
\end{convention}

This ordering is fixed by \code{pack\_complex} / \code{unpack\_complex} and is
assumed by every real-packed matrix assembly in \cref{sec:real-coupled}.

\subsection{Convergence criterion}

The natural scale for the residual is the source that drives it. Define
\begin{equation}
  \boxed{\;
  r_{\mathrm{coeff,rel}}
  =
  \frac{\bigl\|\mathrm{pack}\,\Rvec(\Xvec)\bigr\|_2/\sqrt{N}}
       {\max\!\left(\bigl\|\mathrm{pack}\,\Svec\bigr\|_2/\sqrt{N},\;\epsilon\right)}
  \;<\;
  \varepsilon_{\mathrm{Newton}},
  \;}
  \label{eq:convergence}
\end{equation}
with $N = 2Hn$ for the full problem, or $N = 2Hm$ after Schur reduction to $m$
retained nodes (\cref{sec:schur}). The $1/\sqrt{N}$ normalisation makes the
criterion an RMS-per-unknown quantity, so it does not tighten spuriously as the
circuit grows.

\begin{implnote}[\code{norms}]
\begin{lstlisting}
coeff_abs = norm(R_flat) / sqrt(R_flat.size)
src_abs   = norm(S_flat) / max(sqrt(S_flat.size), 1.0)
coeff_rel = coeff_abs / max(src_abs, 1e-30)
\end{lstlisting}
so $\epsilon = 10^{-30}$. Verified against \eqref{eq:convergence} evaluated
independently: agreement to twelve significant digits.

An optional \code{time\_abs}/\code{time\_rel} pair evaluates the residual in the
time domain instead; it is a diagnostic and is gated off by default (and see
\cref{impl:source-time}).
\end{implnote}

\subsection{Guards}

Before each Newton step the solver checks that the state and its residual are
finite, and it enforces a wall-clock deadline. A non-finite state aborts
immediately rather than propagating \code{NaN} through a factorisation.

\section{Layer 5: residual evaluation}
\label{sec:residual-eval}

Evaluating $\Rvec$ is the alternating frequency--time cycle. Given $\Xvec$:

\paragraph{Step 1 --- synthesise the node-flux waveform.}
\begin{equation}
  \boxed{\;
  \xvec_r = 2\,\Real\!\left[\sum_{k\in\modeset} X_k e^{ik\theta_r}\right],
  \qquad r=0,\dots,N_t-1 .
  \;}
  \label{eq:aft-1}
\end{equation}

\paragraph{Step 2 --- map to branch fluxes.}
\begin{equation}
  \boxed{\;\psi_r = \psi_{\mathrm{dc}} + \Bphi^{\mathsf T}\xvec_r .\;}
  \label{eq:aft-2}
\end{equation}

\paragraph{Step 3 --- apply the branch law pointwise.}
\begin{equation}
  \boxed{\;
  \iJ{}_{,r} = \Ic \hada \sin\!\left(\frac{\psi_r}{\rphi}\right)
  \;-\;
  \Ic \hada \sin\!\left(\frac{\psi_{\mathrm{dc}}}{\rphi}\right) .
  \;}
  \label{eq:aft-3}
\end{equation}
The second term is the static-current subtraction of
\cref{impl:dc-subtraction}. It vanishes identically for
$\psi_{\mathrm{dc}}=\vect{0}$.

\paragraph{Step 4 --- map back to nodes.}
\begin{equation}
  \boxed{\;\vect{n}_r = \Bphi\,\iJ{}_{,r} .\;}
  \label{eq:aft-4}
\end{equation}

\paragraph{Step 5 --- project onto harmonics.}
\begin{equation}
  \boxed{\;
  N_k(\Xvec)
  = \frac{1}{N_t}\sum_{r=0}^{N_t-1}\vect{n}_r\,e^{-ik\theta_r} .
  \;}
  \label{eq:aft-5}
\end{equation}

\paragraph{Step 6 --- form the residual.} Equation~\eqref{eq:residual-k}.

Schematically,
\begin{equation}
  \Xvec
  \xrightarrow{\ \text{synthesis}\ } \xvec(t)
  \xrightarrow{\ \text{branch law}\ } \iJ(t)
  \xrightarrow{\ \text{projection}\ } \Nvec(\Xvec)
  \xrightarrow{\quad} \Rvec(\Xvec).
\end{equation}

The cost is dominated by two sparse products with $\Bphi$ and two dense
$(N_t\times H)$ products, all $O(N_t \cdot \nnz)$ --- cheap compared with a
factorisation.

\section{Layer 6: the Jacobian and its action}
\label{sec:jacobian}

\subsection{The junction tangent}

Differentiating \eqref{eq:aft-3} with respect to $\psi$ gives the time-domain
tangent
\begin{equation}
  \boxed{\;
  \gamma_b(t) = \frac{I_{c,b}}{\rphi}\,
                \cos\!\left(\frac{\psi_b(t)}{\rphi}\right),
  \;}
  \label{eq:gamma-t}
\end{equation}
one value per branch per time sample, frozen at the current Newton iterate. Its
Fourier coefficients are
\begin{equation}
  \ghat_\ell = \frac{1}{N_t}\sum_{r=0}^{N_t-1}\gamma(t_r)\,e^{-i\ell\theta_r},
  \label{eq:gamma-hat}
\end{equation}
and the corresponding \emph{tangent-stiffness harmonics} are
\begin{equation}
  \boxed{\;
  \Khat_\ell = \Bphi\,\diag(\ghat_\ell)\,\Bphi^{\mathsf T}
  \;\in\;\mathbb{C}^{n\times n}.
  \;}
  \label{eq:khat}
\end{equation}

\begin{implnote}[verified]
\code{spectral\_tangent\_state} computes \eqref{eq:gamma-hat} as
\code{np.mean(gamma\_t * phase[:, None], axis=0)} --- identical to the
$1/N_t$ sum --- and \eqref{eq:khat} as
\code{Bphi @ sp.diags(gamma\_hat\_ell) @ BphiT}. Checked against a directly
assembled reference at $\ell=2$: maximum deviation $\num{0.000e+00}$.

The offsets assembled are exactly
$\{k-q\}\cup\{k+q\}$ over all pairs $k,q\in\modeset$ --- the set that
\eqref{eq:jacobian-action} below requires and no more.
\end{implnote}

\subsection{The exact harmonic Jacobian}

Perturbing $\Xvec\to\Xvec+\epsilon\Vvec$ and linearising, the real waveform
perturbation is $2\Real[\sum_q V_q e^{iq\wp t}]$, which contains \emph{both}
$e^{+iq\wp t}$ and $e^{-iq\wp t}$. Multiplying by $\gamma(t)$ and projecting
onto $e^{-ik\wp t}$ therefore picks up $\ghat_{k-q}$ from the first and
$\ghat_{k+q}$ from the second:

\begin{equation}
  \boxed{\;
  (\Jmat\Vvec)_k
  = \Dmat_k V_k
  + \sum_{q\in\modeset}
    \left[\,\Khat_{k-q}\,V_q \;+\; \Khat_{k+q}\,\overline{V_q}\,\right].
  \;}
  \label{eq:jacobian-action}
\end{equation}

The second term is the \emph{conjugate coupling}. It is a direct consequence of
$\Rvec$ not being complex-analytic, and dropping it --- as the conventional
mode-coupled construction does --- is the single most consequential
approximation available in this problem (\cref{sec:precond}).

\subsection{Matrix-free evaluation}

Equation~\eqref{eq:jacobian-action} need never be assembled. The same AFT cycle
that evaluates the residual evaluates its directional derivative:

\begin{align}
  \vect{v}(t_r)      &= 2\,\Real\!\left[\sum_{q\in\modeset} V_q e^{iq\theta_r}\right],
  \label{eq:jvp-1}\\
  \delta\psi(t_r)    &= \Bphi^{\mathsf T}\vect{v}(t_r),
  \label{eq:jvp-2}\\
  \delta\iJ(t_r)     &= \gamma(t_r)\hada\delta\psi(t_r),
  \label{eq:jvp-3}\\
  \delta\vect{n}(t_r)&= \Bphi\,\delta\iJ(t_r),
  \label{eq:jvp-4}\\
  \delta N_k         &= \frac{1}{N_t}\sum_{r=0}^{N_t-1}
                        \delta\vect{n}(t_r)\,e^{-ik\theta_r},
  \label{eq:jvp-5}\\
  (\Jmat\Vvec)_k     &= \Dmat_k V_k + \delta N_k .
  \label{eq:jvp-6}
\end{align}

\begin{proposition}
Equations \eqref{eq:jvp-1}--\eqref{eq:jvp-6} compute exactly
\eqref{eq:jacobian-action}, and equal
\begin{equation}
  \Jmat\Vvec = \left.\frac{\mathrm{d}}{\mathrm{d}\epsilon}
                \Rvec(\Xvec + \epsilon\Vvec)\right|_{\epsilon=0}
\end{equation}
analytically --- not by finite differences.
\end{proposition}
\begin{proof}[Sketch]
The synthesis \eqref{eq:jvp-1} makes both $\pm q$ components explicit; the
pointwise product \eqref{eq:jvp-3} convolves them with the spectrum of $\gamma$;
the projection \eqref{eq:jvp-5} selects harmonic $k$, giving $\ghat_{k-q}$
against the $+q$ component and $\ghat_{k+q}$ against the $-q$ component. The
conjugate term is therefore included automatically, without ever forming
$\Khat_{k+q}$.
\end{proof}

This is the crucial cost property of the method: one Jacobian-vector product
costs one AFT cycle, the same as one residual evaluation, and requires no
storage beyond the frozen $\gamma(t)$.

\begin{implnote}[\code{jvp\_coeffs\_with\_tangent}]
\begin{lstlisting}
JV = [self._linear_blocks[h] @ V[h] for h in range(self.H)]
v_t   = self.grid.synthesize(V)
dpsi  = (self.BphiT @ v_t.T).T
di    = tangent.gamma_t * dpsi
dn    = (self.Bphi @ di.T).T
return JV + self.grid.project_positive(dn)
\end{lstlisting}
Exactly \eqref{eq:jvp-1}--\eqref{eq:jvp-6}. The tangent is passed in rather than
recomputed, so all GMRES iterations within one Newton step share it.
\end{implnote}

\section{Layer 7: preconditioning}
\label{sec:precond}

\subsection{The requirement}

GMRES applied to \eqref{eq:newton} with the matrix-free operator converges
slowly: the harmonic-balance Jacobian of a long chain is badly conditioned. A
preconditioner $\Mmat\approx\Jmat$ is needed for which $\Mmat^{-1}\vect{v}$ is
cheap. Four constructions are available, forming a clear accuracy/cost ladder.

\subsection{Block-diagonal preconditioners}

\paragraph{Linear.} Ignore the nonlinearity entirely:
\begin{equation}
  \Mmat^{\mathrm{linear}}_{kq} = \Dmat_k\,\delta_{kq}.
  \label{eq:precond-linear}
\end{equation}
$H$ independent LU factorisations of the (banded) dynamic blocks. Cheapest,
weakest; adequate only at very low drive.

\paragraph{Mean tangent.} Retain the DC component of the junction tangent:
\begin{equation}
  \bar{\gamma} = \frac{1}{N_t}\sum_r \gamma(t_r),
  \qquad
  \Mmat^{\mathrm{mean}}_{kq}
  = \left[\Dmat_k + \Bphi\,\diag(\bar\gamma)\,\Bphi^{\mathsf T}\right]\delta_{kq}.
  \label{eq:precond-mean}
\end{equation}
Still block-diagonal --- still $H$ small factorisations --- but now aware that
the pump has softened the junctions. This is the default.

\subsection{Mode-coupled preconditioner}

Retaining the $(k-q)$ coupling but dropping the conjugate term gives one large
complex matrix:
\begin{equation}
  \Mmat^{\mathrm{spectral}}_{kq} = \Dmat_k\,\delta_{kq} + \Khat_{k-q}.
  \label{eq:precond-spectral}
\end{equation}
One LU of dimension $Hn$ instead of $H$ of dimension $n$. Substantially
stronger, and correct in the limit where the conjugate coupling is small.

\subsection{Exact real-coupled preconditioner}
\label{sec:real-coupled}

The conjugate coupling is \emph{not} small for stiff DC and mutual-inductor
designs. Retaining it requires working in real coordinates, because the map
$V\mapsto \Khat_{k+q}\overline{V}$ is antilinear.

Write the contribution to output mode $k$ from input mode $q$ as
\begin{equation}
  \mat{L}V_q + \mat{P}\overline{V_q},
  \qquad
  \mat{L} = \Khat_{k-q} + \Dmat_k\delta_{kq},
  \qquad
  \mat{P} = \Khat_{k+q}.
\end{equation}
Splitting $\mat{L} = \mat{L}_r + i\mat{L}_i$, $\mat{P} = \mat{P}_r + i\mat{P}_i$
and $V = V_r + iV_i$, and collecting real and imaginary parts:

\begin{equation}
  \boxed{\;
  \begin{bmatrix} \Real \\ \Imag \end{bmatrix}
  \longleftarrow
  \begin{bmatrix}
    \mat{L}_r + \mat{P}_r & \mat{P}_i - \mat{L}_i \\[2pt]
    \mat{L}_i + \mat{P}_i & \mat{L}_r - \mat{P}_r
  \end{bmatrix}
  \begin{bmatrix} V_r \\ V_i \end{bmatrix}.
  \;}
  \label{eq:real-superblock}
\end{equation}

Assembling \eqref{eq:real-superblock} over all $(k,q)$ gives a real matrix of
dimension $2Hn$ that is the \emph{exact} Jacobian, not an approximation.

\begin{implnote}[\code{\_build\_real\_coupled\_matrix}]
The four quadrants are assembled separately and combined
\code{sp.bmat([[jrr, jri], [jir, jii]])}, matching the \code{pack\_complex}
ordering (all real parts, then all imaginary parts). Because
$\Mmat = \Jmat$ exactly, GMRES converges in approximately one iteration.

The correctness of this assembly is not observable from physics: a dropped sign
would only make the preconditioner slower, and every physics gate would still
pass. It is therefore tested directly, quadrant by quadrant, against
\code{problem.real\_coupled\_matrix} in
\file{tests/test\_fast\_coupled\_assembly.py}.
\end{implnote}

\begin{center}
\begin{tabular}{@{}l l l l@{}}
\toprule
Name & Structure & Factorisations & Typical GMRES iterations \\
\midrule
\code{none}             & --- & 0 & many \\
\code{linear}           & \eqref{eq:precond-linear} & $H\times(n)$ & many \\
\code{mean\_tangent}    & \eqref{eq:precond-mean} & $H\times(n)$ & moderate \\
\code{spectral\_coupled}& \eqref{eq:precond-spectral} & $1\times(Hn)$ & few \\
\code{real\_coupled}    & \eqref{eq:real-superblock} & $1\times(2Hn)$ & $\approx 1$ \\
\bottomrule
\end{tabular}
\end{center}

\section{Layer 7b: factorisation}
\label{sec:factorisation}

Once $\Mmat$ is assembled, a sparse LU with row and column permutations is
computed:
\begin{equation}
  \boxed{\;\mat{P}_r\,\Mmat\,\mat{P}_c = \mat{L}\mat{U}.\;}
  \label{eq:lu}
\end{equation}
Applying $\Mmat^{-1}$ is then two triangular solves and two permutations:
\begin{equation}
  \tilde{\vect{v}} = \mat{P}_r\vect{v},\quad
  \mat{L}\vect{y} = \tilde{\vect{v}},\quad
  \mat{U}\vect{w} = \vect{y},\quad
  \vect{z} = \mat{P}_c\vect{w},
\end{equation}
i.e.
\begin{equation}
  \vect{z} = \Mmat^{-1}\vect{v}
           = \mat{P}_c\mat{U}^{-1}\mat{L}^{-1}\mat{P}_r\vect{v}.
\end{equation}
The expensive part is $\Mmat\to(\mat{L},\mat{U})$; the repeated part
$\vect{v}\mapsto\Mmat^{-1}\vect{v}$ is cheap.

\begin{remark}[two distinct LU factorisations]
\label{rem:two-lus}
In the default Schur-reduced pump solver there are \emph{two} unrelated LU
factorisations, and conflating them is a common source of confusion:
\begin{itemize}
\item $\Dmat_{ee,k} = \mat{L}_{e,k}\mat{U}_{e,k}$, the eliminated linear-node
      block. Constant in $\Xvec$, so factored \emph{once per frequency and
      harmonic}, outside the power loop. See \cref{sec:schur}.
\item $\Mmat(\Xvec) = \mat{L}_M\mat{U}_M$, the Newton preconditioner. Depends
      on the current iterate, so it lives \emph{inside} the Newton loop.
\end{itemize}
\end{remark}

Preconditioner reuse across Newton steps (modified Newton) is available but is
measured to be a net loss; see \cref{sec:precond-reuse}.

\section{Layer 7c: the Krylov solve}
\label{sec:gmres}

\subsection{Preconditioned GMRES}

GMRES solves \eqref{eq:newton} by building a Krylov subspace
\begin{equation}
  \krylov_m(\mat{A},\vect{b})
  = \spanof\{\vect{b},\mat{A}\vect{b},\mat{A}^{2}\vect{b},\dots,
             \mat{A}^{m-1}\vect{b}\}
\end{equation}
and choosing the iterate that minimises the residual over it. After Arnoldi
reduction $\mat{A}\mat{V}_m = \mat{V}_{m+1}\overline{\mat{H}}_m$, the
minimisation becomes a small least-squares problem
\begin{equation}
  \vect{c}_m = \argmin_{\vect{c}}
  \left\|\beta\vect{e}_1 - \overline{\mat{H}}_m\vect{c}\right\|_2 ,
  \qquad \beta = \|\vect{b}\|_2 ,
\end{equation}
solved by Givens rotations at negligible cost.

\subsection{Which side the preconditioner is applied}

There are two placements. \emph{Right} preconditioning solves
\begin{equation}
  \Jmat\Mmat^{-1}\vect{y} = -\Rvec,
  \qquad \Delta\Xvec = \Mmat^{-1}\vect{y},
  \label{eq:right-precond}
\end{equation}
so the Krylov space is $\krylov_m(\Jmat\Mmat^{-1}, -\Rvec)$ and the residual
GMRES minimises is the \emph{true} residual $\|-\Rvec-\Jmat\Delta\Xvec\|$.
\emph{Left} preconditioning solves
\begin{equation}
  \Mmat^{-1}\Jmat\,\Delta\Xvec = -\Mmat^{-1}\Rvec,
  \label{eq:left-precond}
\end{equation}
so the Krylov space is $\krylov_m(\Mmat^{-1}\Jmat, -\Mmat^{-1}\Rvec)$, the
correction comes out directly with no recovery step, and the minimised
quantity is the \emph{preconditioned} residual.

\begin{implnote}[the implementation is left-preconditioned]
\label{impl:left-precond}
The solver calls
\begin{lstlisting}
delta_real, info = gmres_call(A=Aop, b=rhs, M=Mop, rtol=..., atol=...)
Delta = unpack_complex(delta_real, shape)     # used directly
\end{lstlisting}
with \code{rhs = -pack\_complex(R)}. SciPy's \code{gmres} applies \code{M} on
the \emph{left}; the docstring of the installed version (SciPy 1.18.0) states
verbatim ``In this implementation, left preconditioning is used'', and the
source applies \code{psolve} to both the initial residual and each Arnoldi
vector. Consistently, no $\Mmat^{-1}$ is applied to the returned
\code{delta\_real}.

The system actually solved is therefore \eqref{eq:left-precond}, not
\eqref{eq:right-precond}.

\textbf{Consequence.} The convergence tolerance \code{gmres\_rtol} is tested
against $\|\Mmat^{-1}(\vect{b}-\Jmat\Delta\Xvec)\|$ relative to
$\|\Mmat^{-1}\vect{b}\|$, not against the true residual. With
\code{real\_coupled}, where $\Mmat=\Jmat$ exactly, both converge in about one
iteration and the distinction is operationally invisible. With the weaker
block-diagonal preconditioners it is not: a reported GMRES residual is a
statement about the preconditioned system, and should be read as such.
\end{implnote}

\subsection{Mid-solve deadline}

A pathological GMRES call can grind to \code{maxiter} long past any per-point
time budget.

\begin{implnote}[per-iteration abort]
\code{gmres\_call} wraps the user callback and checks elapsed wall time on
\emph{every} GMRES iteration, using \code{callback\_type="pr\_norm"} (which
fires per iteration, not per restart cycle), raising
\code{GmresDeadlineExceeded} from inside the callback. SciPy does not catch
callback exceptions, so this aborts the single call immediately. The previous
scheme, which checked only between Newton iterations, let one call run
$\sim\SI{200}{\second}$ past a \SI{14}{\second} budget.
\end{implnote}

\section{Layer 8: the line search}
\label{sec:line-search}

A full Newton step can overshoot badly when far from the root. The solver
therefore damps the step by backtracking.

Starting from $\alpha_0=1$, generate
\begin{equation}
  \boxed{\;\alpha_j = 2^{-j},\qquad j=0,1,2,\dots\;}
\end{equation}
and for each trial state $\Xvec_{\mathrm{trial},j} = \Xvec + \alpha_j\Delta\Xvec$
recompute the \emph{full nonlinear} residual. Accept the first $\alpha_j$ with
\begin{equation}
  \boxed{\;
  \left\|\Rvec(\Xvec + \alpha_j\Delta\Xvec;\lambda)\right\|
  <
  \left\|\Rvec(\Xvec;\lambda)\right\| ,
  \;}
  \label{eq:armijo}
\end{equation}
and set $\Xvec \leftarrow \Xvec + \alpha_j\Delta\Xvec$. If $\alpha_j$ falls
below $\alpha_{\min}$ without acceptance, the Newton iteration reports a
line-search failure.

\begin{implnote}[simple decrease, in the relative norm]
\begin{lstlisting}
alpha = 1.0
while alpha >= s.min_alpha:
    Xtrial = X + alpha * Delta
    trial_nrm = problem.norms(Xtrial, source_scale, s.compute_time_residual)
    if trial_nrm["coeff_rel"] < nrm["coeff_rel"]:
        accepted = True; break
    alpha *= 0.5
\end{lstlisting}
Two points of detail. First, the acceptance test uses
$r_{\mathrm{coeff,rel}}$ of \eqref{eq:convergence} rather than the bare
$\|\Rvec\|$; since \code{source\_scale} is fixed within one Newton solve, the
two differ by a positive constant and the tests are equivalent. Second, this is
a \emph{simple-decrease} condition, not an Armijo condition with a sufficient-
decrease constant. Simple decrease can in principle stall; in practice, combined
with the exact preconditioner and continuation, it has not been observed to.
\end{implnote}

\section{The complete hierarchy}
\label{sec:hierarchy}

\begin{center}
\begin{minipage}{0.95\textwidth}
\small\ttfamily
\begin{tabbing}
xx\=xx\=xx\=xx\=xxxxxxxxxxxxxxxxxxxxxxxxxxxxxx\=\kill
Frequency loop --- select $f_p$, form the linear operators \\
\> $\omega_p = 2\pi f_p$ \\
\> $D_k(\omega_p) = K - (k\omega_p)^2 C + i k\omega_p G$ \\
\> [Schur:] $\mathcal{S}_k = D_{nn,k} - D_{ne,k}D_{ee,k}^{-1}D_{en,k}$,\quad
   $P_r D_{ee,k} P_c = L_{e,k}U_{e,k}$ \\[4pt]
\> Power loop --- convert source power to pump forcing \\
\>\> $I_p = \sqrt{2P_{\mathrm{device}}/Z_0}$ \\
\>\> $S_k = \tfrac12 I_{\mathrm{solver}}\,e_p\,\delta_{k,k_{\mathrm{src}}}$ \\
\>\> warm start: $X_{\mathrm{pred}} = X_j + \frac{I_{j+1}-I_j}{I_j-I_{j-1}}(X_j-X_{j-1})$ \\[4pt]
\>\> Continuation loop --- scale the source from easy to physical \\
\>\>\> $R(X;\lambda) = \mathcal{D}X + N(X) - \lambda S = 0$,\quad $\lambda: 0\to1$ \\
\>\>\> success: $\Delta\lambda \leftarrow \min(1, 1.5\,\Delta\lambda)$;\quad
       failure: $\Delta\lambda \leftarrow 0.5\,\Delta\lambda$ \\
\>\>\> tangent: $J\dot{X} = S$,\quad $X_{\mathrm{pred}} = X + \Delta\lambda\,\dot{X}$ \\[4pt]
\>\>\> Newton loop --- linearise and seek the root \\
\>\>\>\> $J(X^{(n)})\Delta X^{(n)} = -R(X^{(n)})$ \\
\>\>\>\> $X^{(n+1)} = X^{(n)} + \alpha_n \Delta X^{(n)}$ \\
\>\>\>\> stop when $r_{\mathrm{coeff,rel}} < \varepsilon$ \\[4pt]
\>\>\>\> $\vdash$ Residual: $x_r = 2\Real[\sum_k X_k e^{ik\theta_r}]$;\quad
         $\psi_r = \psi_{\mathrm{dc}} + B_\phi^{\mathsf T}x_r$ \\
\>\>\>\> \phantom{$\vdash$} $i_{J,r} = I_c\odot\sin(\psi_r/\varphi_0) - I_c\odot\sin(\psi_{\mathrm{dc}}/\varphi_0)$ \\
\>\>\>\> \phantom{$\vdash$} $N_k = \frac{1}{N_t}\sum_r B_\phi i_{J,r} e^{-ik\theta_r}$;\quad
         $R_k = D_kX_k + N_k - \lambda S_k$ \\[3pt]
\>\>\>\> $\vdash$ Preconditioner: $\gamma(t) = \frac{I_c}{\varphi_0}\odot\cos(\psi(t)/\varphi_0)$ \\
\>\>\>\> \phantom{$\vdash$} $\widehat{K}_\ell = B_\phi\diag(\hat\gamma_\ell)B_\phi^{\mathsf T}$ \\
\>\>\>\> \phantom{$\vdash$} $(MV)_k = D_kV_k + \sum_q[\widehat{K}_{k-q}V_q + \widehat{K}_{k+q}\overline{V_q}]$ \\[3pt]
\>\>\>\> $\vdash$ Factorisation: $P_r M P_c = LU$;\quad
         $M^{-1}v = P_c U^{-1}L^{-1}P_r v$ \\[3pt]
\>\>\>\> $\vdash$ GMRES (LEFT preconditioned): $M^{-1}J\,\Delta X = -M^{-1}R$ \\
\>\>\>\> \phantom{$\vdash$} $\mathcal{K}_m = \spanof\{b, Ab, \dots, A^{m-1}b\}$,\quad
         $A = M^{-1}J$ \\[3pt]
\>\>\>\> $\vdash$ JVP: $v(t) = 2\Real[\sum_q V_q e^{iq\omega_p t}]$;\quad
         $\delta\psi = B_\phi^{\mathsf T}v$ \\
\>\>\>\> \phantom{$\vdash$} $\delta i_J = \gamma\odot\delta\psi$;\quad
         $(JV)_k = D_kV_k + \mathcal{F}_k[B_\phi\,\delta i_J]$ \\[3pt]
\>\>\>\> $\vdash$ Line search: $\alpha_j = 2^{-j}$;\quad accept if
         $\|R(X+\alpha_j\Delta X)\| < \|R(X)\|$
\end{tabbing}
\end{minipage}
\end{center}

\medskip
The whole numerical process is the chain
\begin{equation}
  f_p \to \Dmat_k
  \;\to\;
  P_p \to \Svec
  \;\to\;
  \lambda \to \Rvec(\Xvec;\lambda)
  \;\to\;
  \Jmat\Delta\Xvec = -\Rvec
  \;\to\;
  \Mmat\to\mat{L}\mat{U}
  \;\to\;
  \text{GMRES}\to\Delta\Xvec
  \;\to\;
  \alpha \to \Xvec_{\mathrm{new}},
\end{equation}
iterated until $\|\Rvec(\Xvec_{\mathrm{new}})\|\to 0$.

\section{Summary of implementation departures}

Three places where the running code differs from the idealised specification
above, collected for reference:

\begin{center}
\begin{tabular}{@{}p{0.20\textwidth} p{0.34\textwidth} p{0.36\textwidth}@{}}
\toprule
\textbf{Item} & \textbf{Specification} & \textbf{Implementation} \\
\midrule
Preconditioner side
& Right: $\Jmat\Mmat^{-1}\vect{y}=-\Rvec$, then $\Delta\Xvec=\Mmat^{-1}\vect{y}$
& \textbf{Left}: $\Mmat^{-1}\Jmat\Delta\Xvec = -\Mmat^{-1}\Rvec$. SciPy
  \code{gmres} preconditions on the left; the returned vector is used directly.
  \code{rtol} therefore measures the preconditioned residual.
  (\cref{impl:left-precond}) \\
\addlinespace
Branch current
& $\iJ = \Ic\hada\sin(\psi/\rphi)$ with $\psi = \psi_{\mathrm{dc}}+\Bphi^{\mathsf T}\xvec$
& The static current $\Ic\hada\sin(\psi_{\mathrm{dc}}/\rphi)$ is
  \textbf{subtracted}. Identically zero for unbiased 4WM. Verified:
  $\num{2.647e-23}$ against the subtracted form, $\num{1.384e-7}$ against the
  literal one. (\cref{impl:dc-subtraction}) \\
\addlinespace
Time-domain source
& $\isrc(t)=2\Real[\Svec_{k_{\mathrm{src}}}e^{ik_{\mathrm{src}}\wp t}]$
& \code{source\_time} hardcodes $\cos(\wp t)$, ignoring $k_{\mathrm{src}}$.
  Deviation $\num{1.536}$ at $k_{\mathrm{src}}=3$. Diagnostic-only path;
  never enters Newton. (\cref{impl:source-time}) \\
\bottomrule
\end{tabular}
\end{center}

Everything else in this chapter --- $\Dmat_k$, the source coefficient, the
residual, the convergence criterion, $\ghat_\ell$, $\Khat_\ell$, the four
preconditioners, the secant and tangent predictors, the adaptive constants
$1.5/0.5$, and the $2^{-j}$ line search --- was verified against the running
code, in several cases to machine precision.
